% 1
\ejercicio{ Si $\mathbb{K}=\mathbb{F}_p$ es un cuerpo de p-elementos:
\begin{enumerate}
    \item ¿Cuantos puntos tiene la recta proyectiva $\mathbb{P}(\mathbb{K}^2)$?
    \item ¿Cuantos puntos tiene el espacio proyectivo de dimensión n sobre $\mathbb{K}$, $\mathbb{P}(\mathbb{K}^{n+1})$?
\end{enumerate}
}
{
    Sea $V$ espacio vectorial sobre el cuerpo $\mathbb{K}$. \\
    \[
    \mathbb{P}(V) = \frac{V \setminus \{0\}}{\sim}
    \]  
    \[
    u \sim v \iff \exists \lambda \in \mathbb{K} \setminus \{0\} : u = \lambda v
    \]
    Para $V = \mathbb{K}^2$ (donde $\mathbb{K} = \mathbb{F}_p$):
    \[
    V \setminus \{0\} = \{(a,b) : a,b \in \mathbb{F}_p \} \setminus \{(0,0)\}
    \]
    Luego:
    \[
    \#(V \setminus \{0\}) = p^2 - 1
    \]
    Ahora tenemos que:
    \[
    (a,b) \sim (c,d) \iff \exists \lambda \in \mathbb{F}_p \setminus \{0\} : (a,b) = (\lambda c, \lambda d)
    \]
    Finalmente calculemos:
    \[
    \#(\mathbb{P}(\mathbb{K}^2)) = \frac{p^2 - 1}{p - 1} = p + 1
    \]
    En general:
    \[
    \#(\mathbb{P}(\mathbb{K}^{n+1})) = \# \left( \frac{\mathbb{K} \setminus \{0\}}{\sim} \right)
    \]
    \[
    \mathbb{K}^{n+1} = (a_0, a_1, \ldots, a_n) : a_i \in \mathbb{K} \implies \#(\mathbb{K}^{n+1} \setminus \{0\}) = p^{n+1} - 1
    \]
    Asi finalmente:
    \[
    \#(\mathbb{P}(\mathbb{K}^{n+1})) = \frac{p^{n+1} - 1}{p - 1} = p^n + p^{n-1} + \ldots + p + 1
    \]
}

% 2
\ejercicio{}{}

% 3
\ejercicio{}{}

% 4
\ejercicio{}{}

% 5
\ejercicio{}{}

% 6
\ejercicio{
Demuestra las siguientes afirmaciones, válidas en un espacio proyectivo cualquiera $\mathbb{P} = \mathbb{P}(V)$, donde $V$ es un espacio vectorial sobre un cuerpo $\mathbb{K}$:
\begin{itemize}
    \item[a)] Por dos puntos distintos pasa una única recta.
    \item[b)] Una recta y un hiperplano que no la contiene se cortan en exactamente un punto.
    \item[c)] Si $\dim \mathbb{P} = 3$, tres hiperplanos distintos se cortan en un punto o una recta.
    \item[d)] Por tres puntos no alineados pasa un plano.
    \item[e)] Dos rectas se cortan si y solo si son coplanarias.
\end{itemize}
}{
    \begin{itemize}
        \item [a)] Sean $A, B \in \mathbb{P}$ dos puntos distintos, veamos que $\exists r \subset \mathbb{P}$ recta tal que $A, B \in r$ y que es única.\\
        \[
        r = V(A, B) = P (L[u, v]) \text{ con } A = [u], B = [v] 
        \]
        Como $P(L[u, v])$ tiene dimensión 1, es una recta, y claramente esta contenido en toda recta que contenga a $A$ y $B$, la cual también tiene que ser de dimensión 1, luego es única.
        \item [b)] Sean la recta $r$ y el hiperplano $H$ que no la contiene.\\
        Sabemos que sus dimensiones son:
        \begin{align*}
            \dim(r) &= 1 \\
            \dim(H) &= \dim(\mathbb{P}) - 1
        \end{align*}
        Veamos entonces que $\dim(r \cap H) = 0$, usando que $r \not\subset H$.
        \\
        \[
        \dim(V(r, H)) = \dim(r) + \dim(H) - \dim(r \cap H) \implies
        \]
        \[
        \dim(r \cap H) = \dim(r) + \dim(H) - \dim(V(r, H)) = 1 + (\dim(\mathbb{P}) - 1) - \dim(\mathbb{P}) = 0
        \]
        Esto se da ya que:
        \[
        r, H \subset V(r, H) \implies \dim(V(r, H)) \geq \dim(H) = \dim(\mathbb{P}) - 1 \\
        \]
        Como $r \not\subset H$, entonces $\dim(V(r, H)) > \dim(H)$, ya que si no serian dos hiperplanos iguales al compartir la misma dimensión y contenerse uno a otro, ya que claramente $H \subset V(r, H)$, luego $\dim(V(r, H)) = \dim(\mathbb{P})$.\\
        Luego $\dim(r \cap H) = 0$, luego $r \cap H$ es un punto.
        \item [c)] Sea $\dim(\mathbb{P}) = 3$ y $H_1, H_2, H_3$ tres hiperplanos distintos.\\
        Tenemos que comprobar que $\dim(H_1 \cap H_2 \cap H_3) \in \{0, 1\}$.\\
        Sea $X = H_1 \cap H_2$ y por la formula de Grassmann:
        \[
        \dim(H_1 \cap H_2 \cap H_3) = \dim(X \cap H_3) = \underbrace{\dim(H_1 \cap H_2)}_{=\dim(X)} + \dim(H_3) - \dim(V(X, H_3))
        \]
        \[
        \dim(X) = \dim(H_1 \cap H_2) = \dim(H_1) + \dim(H_2) - \dim(V(H_1, H_2)) = 2 + 2 - 3 = 1
        \]
        Entonces como $\dim(V(X, H_3)) \geq \dim(H_3) = 2$, tenemos que:
        \[
        \dim(H_1 \cap H_2 \cap H_3) = \dim(X \cap H_3) = \underbrace{\dim(H_1 \cap H_2)}_{=\dim(X)} + \dim(H_3) - \dim(V(X, H_3)) \leq 1 + 2 - 2 = 1
        \]
        Ademas $\dim(H_1 \cap H_2 \cap H_3) \neq -1$ ya que si no:
        \[
        \dim(H_1 \cap H_2 \cap H_3) = -1 \implies -1 = 1 + 2 - \dim(V(X, H_3)) \implies \dim(V(X, H_3)) = 4
        \]
        Lo cual es claramente absurdo, ya que $\dim(\mathbb{P}) = 3$.\\
        Luego $\dim(H_1 \cap H_2 \cap H_3) \in \{0, 1\}$.
    \end{itemize}
}
